% claim-ref: eisenia-fetida#bin-bedding
% claim-ref: eisenia-fetida#temperature
% claim-ref: eisenia-fetida#startup-feed
% claim-ref: eisenia-fetida#readiness-record
\CompanionHeader{BIN SETUP / NUMBERS \& PACIFICA}{An aerated bedding bin}{Red wiggler / Eisenia fetida / planning only}{Prepare and measure the complete terrestrial habitat before obtaining stock.}
\Context{Home-bin target: 15--25 C (59--77 F), from San Mateo County guidance. Measure the bedding: a Pacifica address, shade and cool air do not prove a safe bin. [EF-S1; proposed]}
\Panel{Prepare the container}{Clean opaque plastic or suitable untreated wood; lid excludes light without sealing out air. Upper screened vents, bottom drainage and separate catch tray; secure outdoors against rodents. About 30 cm (12 in) deep is a shallow design example, not a minimum-volume rule. [EF-EPA; EF-O1; proposed]}
\Panel{Build the bedding}{Start with 10--15 cm (4--6 in) moistened, fluffed bedding. Use plain uncoated shredded cardboard and non-glossy newspaper; exclude pesticides, unknown outdoor organisms and fresh grass that may heat. Keep vents open; squeeze only a few drops. Replenish bedding as consumed. [EF-O1; EF-N1; EF-U1; conservative adaptation]}
\Panel{Place, measure, adjust}{Shelter from direct sun, rain and excessive cold; retain airflow. Record morning and warmest-period bedding temperatures before stocking and after weather changes. Below about 10 C (50 F), feeding slows: that is not a target. Move or insulate before hazardous cold; move to cooler shelter in heat. No thermal or drying stress to force breeding. [EF-S1; EF-O1; proposed]}
\Panel{Start small; observe consumption}{OSU startup example: 0.11--0.23 kg scraps per 0.45 kg worms per WEEK (1/4--1/2 lb per 1 lb). Not a daily ration or promised throughput. Let identified grower stock settle into prepared bedding before small feedings; add more only as consumed. No stock count or capacity is established here. [EF-O1; EF-N1; EF-EPA]}
\Panel{Readiness record / all values unmeasured}{Dimensions / volume / bedding materials and depth: \hrulefill\par
Site / shade / ventilation / drainage / contingency: \hrulefill\par
Source / date / supplier identification / stock estimate / stages: \hrulefill\par
Bedding min/max temperature / squeeze test: \hrulefill\par
Feed offered / remaining / odor / pests / escapes: \hrulefill}
\Note{Before worms arrive}{Verify the whole bin and intended site; obtain traceable composting stock, not unidentified yard or bait worms. Check requirements before import/transfer. Keep an accessible, movable bin and a suitable contingency. [EF-EPA; EF-N1; proposed]}
\Sources{Clickable evidence: \href{https://www.epa.gov/recycle/composting-home}{EF-EPA}; \href{https://extension.oregonstate.edu/catalog/em-9034-composting-worms}{EF-O1}; \href{https://content.ces.ncsu.edu/worms-can-recycle-your-garbage}{EF-N1}; \href{https://ucanr.edu/node/161011/printable/print}{EF-U1}; \href{https://www.smcsustainability.org/waste-reduction/compost/composting-at-home/}{EF-S1}. Claims, source scope, separate dates and contradictions: adjacent YAML. Planning only; habitat unmeasured.}{eisenia-fetida / tank-setup}
